% main.tex - Project: Network Intrusion Detection System
\documentclass[12pt,a4paper]{article}
\usepackage[utf8]{inputenc}
\usepackage[T1]{fontenc}
\usepackage[english]{babel}
\usepackage{csquotes}
\usepackage{microtype}
\usepackage{geometry}
\usepackage{comment}
\geometry{margin=2.5cm}
\setlength{\parindent}{0pt}
\usepackage{setspace}
\usepackage{comment}

\usepackage{minted}
\usepackage{xcolor} 
\definecolor{mybgcolor}{HTML}{000000}

\onehalfspacing

\usepackage{hyperref}
\hypersetup{
    pdftitle={},
    pdfauthor={First Defenders},
    colorlinks=true,
    linkcolor=blue,
    citecolor=blue,
    urlcolor=blue
}

\usepackage{graphicx}
\graphicspath{ {./images/} }
\usepackage{float}
\usepackage{caption}
\usepackage{booktabs}
\usepackage{enumitem}

% Title metadata
\usepackage[utf8]{inputenc}
\usepackage{geometry}
\geometry{top=3cm,bottom=3cm,left=3cm,right=3cm}

\usepackage{setspace}

\begin{document}
\begin{titlepage}

    \centering
    \vspace*{2cm}
    
    {\Huge \makebox[\textwidth][c]{\textbf{Lorem ipsum dolor sit amet}}}\\[2cm]
    
    {\Large P1 Procesanalyse DC1 Grp. 6}\\[1cm]
    
    \includegraphics[width=0.4\textwidth]{report/images/aau_logo.png}
    
    \begin{flushleft}
    \textbf{Authors:} \\
    \hspace*{2cm} Mads Lindegaard\\
    \hspace*{2cm} Emil Malik Stokholm\\
    \hspace*{2cm} Tobias Scholtz\\
    \hspace*{2cm} Dharshika Christopher\\[1cm]
    
    \textbf{Supervisor:} \\
    \hspace*{2cm} Reza Tadayoni\\
    \hspace*{2cm} Mads Bisgaard Schmidt\\[1cm]
    
    \textbf{Semester:} \\
    \hspace*{2cm} 1st Semester\\[0.5cm]
    
    \textbf{Study Programme:} \\
    \hspace*{2cm} Cybersecurity\\[0.5cm]
    
    \end{flushleft}
    
    \vfill
    \textbf{Date:} December 11, 2026
\end{titlepage}

\clearpage

% -------------------------
% 0. Table of Contents
% -------------------------

\clearpage
\tableofcontents
\clearpage

% -------------------------
% 1. Introduction
% -------------------------
\section{Introduktion}
\label{introduktion}
Denne processanalyse har til formål, at belyse vores arbejdsprocess og evaluere den med henblik på fremtidige forbedringer. \\ \\
Vi dannede vores gruppe i begyndelsen af oktober efter et fællesmøde om P1-projektet. Vi ønsket alle at vi vælge et emne som kunne være aktuelt på arbejdsmarkedet, og denne fællesnævneren var begrundelsen for netop vi havnet på gruppe sammen. 


% -------------------------
% 2. Problem Field

% -------------------------

\section{Projektstyring og planlægning}
\subsection{Planlægning og struktur}
Overordnet har vi startet projektet med at kortlægge den langsigtede plan for projektet. Dette vi har udfyldt i et Gant diagram. Vi har fastslået datoer og deadlines, ved at arbejde baglæns fra den endelige eksamensdato og frem til starten af projektet. Hertil har vi så vidt muligt taget højde for weekender og perioder med andre eksamensdatoer.

Vi har benyttet en fælles kalender med midlertidige møde datoer, for at sikre os som minimum et til to ugentlige møde igennem hele projektet. Møderne har som udgangspunkt i starten af projektet, haft til formål at vidensdele blandt gruppemedlemerne, samt koordinere og udeligere opgaver. Senere hen i forløbet er der blevet tilføjet yderligere dagsordner, som bla. vejledermøde, forberedelse og fælles skrivetidspunkter, hvor vi har siddet i gruppen og skrevet individuelt, for herefter at gennemgå i planum.

\\
 I starten af processen har vi fokuseret på at få udbygget vores generelle viden om emnet, og problematikker indenfor emnet for at kunne lave en midlertidig problemformulering. Denne fase har vi kategoriseret i vores overordnede plan som undersøgelsesfase. Denne fase har været defineret af bred vidensøgning, og høj grad af vidensdelling, samt diskution omkring emnet til vores møder. 
 Undersøgelsesfasen har konkret primært bestået af kortsigtede delmål, som der er blevet aftalt og udeligeret på de enkelte møder. Dette har fungeret for at vi ikke skuller overlappe hinanden i undersøgelsesområder, og for at sikre at alle gruppemedlemer individuelt har haft relevante arbejdsopgaver. Vi har efterfølgende fra undersøgelsesfasen revurderet og tiljusteret vores overordnede tidsplan, til at deadlinsne mere realistisk afspejler det egentlige tidsforbrug retrospektivt, og fremadrettet, for at præcisere mere realistiske tidsestimeringer til fremtidige opgaver og perioder.  
\\
som vi er gået fra undersøgelsesfase over til skrive og produktionsfase, har vores mødedagsordner ændret karakter så de passer bedre til denne process. Vi har i høj grad skåret ned på informationsdelling, for at fokusere mere på brainstorming, samt forberedelse på vejledermøder og generelle skrivesessioner. Dette har ikke været relevant tidligere i undersøgelsesfasen, da vores vejlednings møder indtil denne fase i høj grad har bestået at teknske og opgaveformat mæssige spørgsmål, samt generelle spørgsmål til planlægning. Hvorimod vejledninger i skrive og produktionsfasen i høj grad har taget udgangspunkt i materialet der er blevet produceret fra gang til gang, og benyttet til at sætte nye delmål samt retning på skrivning og produkt format. 
\\
\\


Måden vi samarbejder på: videosamtaler ved sygdom, 
gruppemøder, forventningsafstemning i forbindelse med hvad vi bedst kan lide (alle kan bedst lide at mødes)


\subsection{Organisering, roller og opgavefordeling}
opgavefordeling har vi som gruppe besluttet os at prøve at arbejde med kanbas diagram for at fordele og orginasier opgaver mellem os. Her benyttet vi os af Loop microsoft applikationer. Grunden til vi valgte loop var vi gerne vil have at vi ikke brugte for mange forskellige applikationer men istedet for prøve at samle det en sted så mest muligt. Da vi alle have problemmer med i vores tideligere projekt, at det var svært at få overblik over hvad ligger hvor og hvad anvender vi til hvad. Dermed endte med at miste motivationen til at orginiserer opgave fordeling. Hermed tænkte vi hvis vi brugte microsft applikationer for at minimere antal af applikationer, som taler sammen, vil det give bedre overblik og motivation til benytte kanban diagramet. Måden vi organiseret fordelling af opgaver på var at vi fordeler opgaver er kun til møder og det hvad vi slutter med som til mødet. Her skriver vi ned de opgaver der skal laves, derefter Så spørg vi hinanden hvad folk synes er spænenden at tage sig til, hvis nogle der vil arbejde med det samme så laver vi....

dermed når opgaverne er fordelt sættes en fælles deadline til, hvornår opgaverne forventes at være færdig.  

indtil videre har vi konkluderet at vi anvender det mere med kanbat end hvad vi førhen gjorte men stadig føler at vi mister overblik, da vi enig om vi stadig har for mange applikationer vi anvender. vi har diskuteret vi vil prøve at anvende den gamledags metode med stickers på en tavle i vores projektrum. 

roller:
Udfra de møder vi har haft løbet af projektet, har vi kunne begynde at antyde hvilken roller hinanden hælder mest til. 
vi har idemageren, kordinatioren, 
\clearpage

\section{Problemanalyse}
\subsection{Problemformulering}
For at finde frem til en relevant problemformulering, har vi benyttet opgaveteksten fra vores emnevalg, Sikker filopbevaring og –deling, til at insnævre det initerende problem, som vores problemformulering tager udgangspunkt i. Vores midlertidige problemformulering er på nuværerende tidspunkt: \\
Which solutions can improve the security in relation to file sharing in small organizations?\\


\subsection{Vejledersamarbejde}
Vi startede vores samarbejde ved at planlægge tidspunkter hvor vi skulle mødes. Vores samtaler med vejleder blev planlagt efter vores allerede aftale møder. 

Til nu har vi kun haft to samtaler med vejleder, men vi skal fra skrivende stund sende et udkast til vejleder hver uge, selv om vi ikke kan mødes hver uge. \\
1. møde med vejleder: struktur \\
2. møde med vejleder: blev anbefalet at programmere i C og ikke lave en server. 

\subsection{Brug af AI}
Vi har brugt AI på forskellige måder tidligere. Vi blev enige om at vi kan bruge søgemotorer som er koblet til AI (eksempelvis i Primo AAU og google), men ikke at bruge det til retsskrivning og kode. 
\clearpage

\section{Samarbejde}
\subsection{Akademisk}
Vi er alle over 25 og har studeret tidligere, så vi havde en forventning om en viss akademisk standard. Dette gælder både i form af kildehåndtering og referencer, så vel som AI.\\ \\
Vores akademiske roller blev fordelt efter individuelle kompetencer. Dharshika er meget komfortabel på engelsk, og havde derfor hovedansvaret for grammatiske rettelser og projektets underliggende tone. Med sin baggrund som datamatiker, var Tobias ansvarlig for de større beslutninger indenfor kodning af produktet. Emil har studeret ved AAU Aalborg og delte sin erfaring med gruppearbejde og brainstorming. Mads var altid forud med sit arbejde og overblik og fungerede som gruppens initiativtager. 



\subsection{Socialt}
Da vi dannede gruppen, var vi alle enige om at et vigtigt element for et godt samarbejde var et socialt fællesskab i gruppen. Efter nogle møder indførede vi et socialt kvarter, nemlig at vi i de første 15 minutter af et aftalt møde kunne tale helt frit om vores dag, humør og andet vi havde på hjertet. Når vi har haft møder i længere perioder, har vi også haft aktive pauser, hvor vi har gået udenfor og spillet basketbold.  

\subsection{Konflikthåndtering}
Ikke rigtigt en konflikt, men vi har haft sygdom i gruppen. Dette løste vi ved at vedkommende kunne deltage i møderne online over Discord.  


% ---------------------
% 11. Appendix (optional)
% ---------------------
\section{Evaluering}
\subsection{Hvad har vi lært?}
\subsection{Reflektioner}
\subsection{Forbedringspunkter}
\subsection{Hvad kan vi tages med til P2?}

\clearpage
% ---------------------
% References
% ---------------------
\section{Bilag}


\end{document}